\section{部分分数}\label{節:部分分数}

割り算を終えても,分母が2つの一次式の積になっていることがある.一般的な形は
\[
 \int\frac{px+q}{(x-\alpha)(x-\beta)}\dd x\qquad(\alpha\ne\beta)
\]
で与えられる.分数の分母を勝手に分配することはできない.たとえば $\frac1{x(x+1)}$ は $\frac1x+\frac1{x+1}$ ではない.しかし,一次式の逆数ならば対数で積分できる.そこで,この分数を $\frac A{x-\alpha}$ と $\frac B{x-\beta}$ の和に書き換えて,対数に帰着させることを目標とする.この書き換えを\textbf{部分分数分解}という.

ここで1つ断りをいれておこう.本節の分数は,分子の次数が分母より低いもの,すなわち割り算を済ませたものとする.また,分母を $0$ にする $x=\alpha,\beta$ では定義されないので,答えはそれらの点をまたがない各区間で読む.

右辺を通分すると
\[
 \frac A{x-\alpha}+\frac B{x-\beta}
 =\frac{A(x-\beta)+B(x-\alpha)}{(x-\alpha)(x-\beta)}
\]
となる.したがって,分子どうしの等式
\[
 px+q=A(x-\beta)+B(x-\alpha)
\]
が $x$ についての恒等式となるように $A,B$ を選べばよい.\,$x=\alpha$ を代入すると $B$ の項が消えて
\[
 p\alpha+q=A(\alpha-\beta)
\]
となり,
\[
 A=\frac{p\alpha+q}{\alpha-\beta}
\]
を得る.同様に $x=\beta$ を代入して
\[
 B=\frac{p\beta+q}{\beta-\alpha}
\]
を得る.これは,元の分数に分母を $0$ にする値を代入したのではない.分母を払った整式の恒等式に代入したのである.逆にこの $A,B$ を選べば,両辺の差は1次以下の整式であり,\,$x=\alpha,\beta$ の2点で $0$ になるから,恒等的に $0$ である.よって
\[
 \int\frac{px+q}{(x-\alpha)(x-\beta)}\dd x
 =A\log|x-\alpha|+B\log|x-\beta|+C
\]
となる.

分母に同じ一次因子が重なって $(x-\alpha)^2$ となっている場合は,上の式で $\beta=\alpha$ とすることはできない.\,$A,B$ の分母が $0$ になるからである.この場合は,分子に分母と同じ $x-\alpha$ を作ってやればよい.
\[
 px+q=p(x-\alpha)+(p\alpha+q)
\]
と書けば,
\[
 \frac{px+q}{(x-\alpha)^2}
 =\frac p{x-\alpha}+\frac{p\alpha+q}{(x-\alpha)^2}
\]
となる.第1項は対数で,第2項は $(x-\alpha)^{-2}$ なのでべきの公式で積分でき,
\[
 \int\frac{px+q}{(x-\alpha)^2}\dd x
 =p\log|x-\alpha|-\frac{p\alpha+q}{x-\alpha}+C
\]
を得る.

\begin{bookpoint}
\textbf{POINT}\par
分母の因数ごとに $\frac A{x-\alpha}$ の形へ分け,係数を恒等式で決めて対数で積分する.重なる因子 $(x-\alpha)^2$ は,分子を $x-\alpha$ と定数の和に書いて分ける.
\end{bookpoint}

\begin{example}{部分分数}
$\int\frac{5x+1}{(x-1)(x+2)}\dd x$ を求めよ.
\end{example}

\paragraphbf{方針}
$\frac A{x-1}+\frac B{x+2}$ と置き,分母を払った恒等式に $x=1,-2$ を代入して係数を決める.

\paragraphbf{解答}
$2\log|x-1|+3\log|x+2|+C$

\paragraphbf{解法}
\[
 \frac{5x+1}{(x-1)(x+2)}
 =\frac A{x-1}+\frac B{x+2}
\]
と置き,分母を払うと
\[
 5x+1=A(x+2)+B(x-1)
\]
となる.\,$x=1$ を代入すると $6=3A$,\,$x=-2$ を代入すると $-9=-3B$ なので,\,$A=2$,\,$B=3$ である.したがって
\[
 \int\frac{5x+1}{(x-1)(x+2)}\dd x
 =2\log|x-1|+3\log|x+2|+C
\]
を得る.元の関数は $x=-2,1$ で定義されないので,この答えは $x<-2$,\,$-2<x<1$,\,$1<x$ の各区間で成り立つ.

\begin{example}{重複する因子}
$\int\frac{2x+3}{(x+1)^2}\dd x$ を求めよ.
\end{example}

\paragraphbf{方針}
分子を $2(x+1)+1$ と書き,\,$\frac2{x+1}$ と $\frac1{(x+1)^2}$ に分ける.

\paragraphbf{解答}
$2\log|x+1|-\frac1{x+1}+C$

\paragraphbf{解法}
分母にある $x+1$ を分子にも作りたいので,\,$2x+3=2(x+1)+1$ と書く.すると
\[
 \frac{2x+3}{(x+1)^2}
 =\frac2{x+1}+\frac1{(x+1)^2}
\]
となる.第1項には対数,第2項にはべきの公式を使い,
\[
 \int\frac{2x+3}{(x+1)^2}\dd x
 =2\log|x+1|-\frac1{x+1}+C
\]
を得る.\,$x=-1$ を除き,\,$x<-1$ と $x>-1$ の各区間で成り立つ.

\paragraphbf{係数比較による決め方}
係数は,恒等式の係数を比べて決めることもできる.本節の1つ目の例題では,右辺を整理すると
\[
 5x+1=(A+B)x+(2A-B)
\]
となるから,
\[
 A+B=5,\qquad 2A-B=1
\]
を解いて,同じ $A=2$,\,$B=3$ を得る.係数比較では連立方程式を解くことになるが,\,$x=\alpha$ を代入する方法では,1回の代入で係数が1つずつ決まる.分母の因数が3つ以上に増えても,この代入は同じように使える.

一般に有理関数,すなわち整式どうしの商を積分するときは,まず割り算で分子の次数を下げ,次に分母を因数分解して分けることを考える.ただし,実数の範囲で一次式に分かれない二次式もある.その場合に役立つ形を次節で調べる.
